% !TEX program = xelatex
% Самодостаточный конспект. Компилировать XeLaTeX/LuaLaTeX (tectonic 06-chislovye-ryady.tex).
\documentclass[11pt,a4paper]{article}
\newcommand{\coursename}{Математический анализ}
\newcommand{\rightheader}{§6. Числовые ряды}
\newcommand{\doctitle}{Числовые ряды}
% ==== Общая преамбула конспектов (собирается в каждый .tex как самодостаточная) ====
% Перед этим файлом должны быть определены: \documentclass и макросы
%   \coursename  — название курса (левый колонтитул, подзаголовок),
%   \rightheader — правый колонтитул (напр. «§2. Рациональные дроби»),
%   \doctitle    — заголовок раздела на титуле.
\usepackage{fontspec}
\setmainfont{cmunrm.otf}[
  BoldFont       = cmunbx.otf,
  ItalicFont     = cmunti.otf,
  BoldItalicFont = cmunbi.otf]
\usepackage{polyglossia}
\setdefaultlanguage{russian}
\setotherlanguage{english}

\usepackage{amsmath,amssymb,amsthm,mathtools}
\usepackage[a4paper,margin=2.2cm]{geometry}
\usepackage{enumitem}
\usepackage{graphicx}
\usepackage{array}
\usepackage{xcolor}
\definecolor{accent}{HTML}{1F6FEB}
\definecolor{rulecol}{HTML}{D0D7DE}
\usepackage[colorlinks=true,linkcolor=accent,urlcolor=accent,citecolor=accent]{hyperref}
\usepackage{titlesec}
\usepackage{fancyhdr}

% ----- Колонтитулы -----
\pagestyle{fancy}
\fancyhf{}
\fancyhead[L]{\small\itshape \coursename}
\fancyhead[R]{\small\itshape \rightheader}
\fancyfoot[C]{\thepage}
\renewcommand{\headrulewidth}{0.4pt}
\renewcommand{\headrule}{\hbox to\headwidth{\color{rulecol}\leaders\hrule height \headrulewidth\hfill}}

% ----- Заголовки -----
\titleformat{\section}{\large\bfseries\color{accent}}{\thesection.}{0.6em}{}
\titleformat{\subsection}{\bfseries}{\thesubsection.}{0.5em}{}

% ----- Теоремные окружения (стиль конспекта) -----
\theoremstyle{definition}
\newtheorem{definition}{Определение}[section]
\newtheorem{example}[definition]{Пример}
\newtheorem{remark}[definition]{Замечание}
\theoremstyle{plain}
\newtheorem{theorem}[definition]{Теорема}
\newtheorem{lemma}[definition]{Лемма}
\newtheorem{corollary}[definition]{Следствие}
\newtheorem{property}[definition]{Свойство}
\newtheorem{criterion}[definition]{Критерий}
\renewcommand{\proofname}{Доказательство}

% ----- Операторы и сокращения -----
\DeclareMathOperator{\tg}{tg}
\DeclareMathOperator{\ctg}{ctg}
\DeclareMathOperator{\arctg}{arctg}
\DeclareMathOperator{\arcctg}{arcctg}
\DeclareMathOperator{\sh}{sh}
\DeclareMathOperator{\ch}{ch}
\DeclareMathOperator{\Th}{th}
\DeclareMathOperator{\cth}{cth}
\DeclareMathOperator{\Const}{const}
\DeclareMathOperator{\sign}{sign}
\DeclareMathOperator{\rang}{rang}
\DeclareMathOperator{\diam}{diam}
\DeclareMathOperator{\Span}{span}
\DeclareMathOperator{\Ker}{Ker}
\DeclareMathOperator{\Img}{Im}
\DeclareMathOperator{\tr}{tr}
\newcommand{\R}{\mathbb{R}}
\newcommand{\C}{\mathbb{C}}
\newcommand{\N}{\mathbb{N}}
\newcommand{\Z}{\mathbb{Z}}
\newcommand{\Q}{\mathbb{Q}}
\newcommand{\dx}{\,dx}
\newcommand{\dt}{\,dt}
\newcommand{\du}{\,du}
\newcommand{\eps}{\varepsilon}
\renewcommand{\Re}{\operatorname{Re}}
\renewcommand{\Im}{\operatorname{Im}}
\renewcommand{\le}{\leqslant}
\renewcommand{\ge}{\geqslant}
\newcommand{\scal}[2]{\left( #1,\,#2\right)}
\DeclarePairedDelimiter{\abs}{\lvert}{\rvert}
\DeclarePairedDelimiter{\norm}{\lVert}{\rVert}

\begin{document}
\begin{center}
  {\LARGE\bfseries \doctitle}\\[4pt]
  {\large\itshape \coursename}
\end{center}
\vspace{6pt}

\section{Числовой ряд и его сумма}

\begin{definition}
Пусть $\{a_n\}_{n=1}^\infty$ — числовая последовательность. Выражение
$\displaystyle\sum_{n=1}^{\infty}a_n=a_1+a_2+\dots+a_n+\dots$ называют \emph{числовым рядом}, $a_n$
— его \emph{общим членом}. Число $\displaystyle S_n=\sum_{k=1}^{n}a_k$ называют $n$-й
\emph{частичной суммой}.
\end{definition}

\begin{definition}
Ряд \emph{сходится}, если существует конечный предел $\displaystyle S=\lim_{n\to\infty}S_n$,
называемый \emph{суммой ряда}. Если предела нет (или он бесконечен), ряд \emph{расходится}.
\end{definition}

\begin{example}
Ряд $\displaystyle\sum_{k=0}^{\infty}(-1)^k$ расходится: $S_0=1,\ S_1=0,\ S_2=1,\dots$ — предела
нет. (Формальные манипуляции вида «$S=1-S\Rightarrow S=\tfrac12$» здесь незаконны: ряд не
сходится.)
\end{example}

\section{Необходимое условие сходимости}

\begin{theorem}
Если ряд $\displaystyle\sum_{k=1}^{\infty}a_k$ сходится, то $\displaystyle\lim_{k\to\infty}a_k=0$.
\end{theorem}

\begin{proof}
Пусть $S_n\to S$; тогда и $S_{n+1}\to S$. Значит $a_{n+1}=S_{n+1}-S_n\to S-S=0$.
\end{proof}

\begin{corollary}[достаточный признак расходимости]
Если $\displaystyle\lim_{k\to\infty}a_k\ne0$ (или предела нет), то ряд расходится.
\end{corollary}

\section{Эталонные ряды}

\subsection{Геометрический ряд}

\begin{theorem}
Ряд $\displaystyle\sum_{k=0}^{\infty}a\,q^k$ сходится при $\abs q<1$ (с суммой
$\dfrac{a}{1-q}$) и расходится при $\abs q\ge1$.
\end{theorem}

\begin{proof}
При $q\ne1$ имеем $S_n=a\dfrac{1-q^{n+1}}{1-q}$. Если $\abs q<1$, то $q^{n+1}\to0$ и
$S_n\to\dfrac{a}{1-q}$; если $\abs q>1$, то $q^{n+1}\to\infty$. При $q=1$ ряд равен $a+a+\dots$, при
$q=-1$ суммы колеблются — в обоих случаях расходимость.
\end{proof}

\subsection{Гармонический ряд}

\begin{theorem}
Гармонический ряд $\displaystyle\sum_{k=1}^{\infty}\frac1k$ расходится.
\end{theorem}

\begin{proof}
Группируем члены:
\[
  1+\frac12+\Bigl(\frac13+\frac14\Bigr)+\Bigl(\frac15+\dots+\frac18\Bigr)+\dots
   \ge 1+\frac12+\frac12+\frac12+\dots=+\infty,
\]
так как в группе из $2^{m}$ слагаемых каждое $\ge\dfrac1{2^{m+1}}$, и сумма группы $\ge\tfrac12$.
\end{proof}

\begin{remark}[скорость расходимости]
$\displaystyle H_n=\sum_{k=1}^n\frac1k=\ln n+\gamma+\eps_n$, где $\eps_n=O(1/n)$, а
$\gamma=\displaystyle\lim_{n\to\infty}(H_n-\ln n)\approx0{,}5772$ — \emph{постоянная
Эйлера–Маскерони}. Существование предела доказывается монотонностью и ограниченностью
$\gamma_n=H_n-\ln n$ (сравнение суммы с интегралом $\int_1^n\frac{dx}{x}=\ln n$).
\end{remark}

\subsection{Телескопический ряд}

\begin{example}
$\displaystyle\sum_{k=1}^\infty\frac{1}{k(k+1)}$: поскольку $\dfrac{1}{k(k+1)}=\dfrac1k-\dfrac1{k+1}$,
\[
  S_n=\Bigl(1-\tfrac12\Bigr)+\Bigl(\tfrac12-\tfrac13\Bigr)+\dots+\Bigl(\tfrac1n-\tfrac1{n+1}\Bigr)
   =1-\frac1{n+1}\to1 .
\]
\end{example}

\begin{remark}[знаменитый ряд Эйлера]
$\displaystyle\sum_{k=1}^\infty\frac1{k^2}=\frac{\pi^2}{6}$ (получается из разложения
$\dfrac{\sin x}{x}=\prod_{k\ge1}\bigl(1-\tfrac{x^2}{\pi^2k^2}\bigr)$ сравнением коэффициентов при
$x^2$).
\end{remark}

\section{Критерий Коши}

\begin{criterion}[Коши]
Ряд $\displaystyle\sum_{k=1}^\infty a_k$ сходится тогда и только тогда, когда
\[
  \forall\eps>0\ \exists N\ \forall n>N\ \forall p\in\N:\quad
  \abs{S_{n+p}-S_n}=\abs[\Big]{\sum_{k=n+1}^{n+p}a_k}<\eps .
\]
\end{criterion}

\begin{proof}
Это критерий Коши сходимости последовательности $\{S_n\}$: $\{S_n\}$ сходится $\iff$ фундаментальна,
а $S_{n+p}-S_n=a_{n+1}+\dots+a_{n+p}$.
\end{proof}

\begin{example}
Расходимость $\sum1/k$ повторно: при $p=n$
$\displaystyle\abs{S_{2n}-S_n}=\frac1{n+1}+\dots+\frac1{2n}\ge\frac1{2n}\cdot n=\frac12$, так что
условие Коши не выполнено.
\end{example}

\section{Свойства сходящихся рядов}

\begin{property}[линейность]
Если $\sum a_k=A$ и $\sum b_k=B$ сходятся, то сходятся $\sum(a_k+b_k)=A+B$ и
$\sum\lambda a_k=\lambda A$ ($\lambda\in\R$).
\end{property}

\begin{property}
Отбрасывание или добавление конечного числа членов не влияет на сходимость (меняет лишь сумму).
\end{property}

\begin{proof}
Если отбросить первые $K$ членов, новая частичная сумма $\sigma_n$ связана со старой соотношением
$S_n=\sigma_n+C$, где $C$ — сумма отброшенных членов. Постоянное слагаемое не влияет на наличие
предела.
\end{proof}

\section{Ряды с неотрицательными членами}

Для $a_k\ge0$ частичные суммы монотонно не убывают, поэтому ряд сходится тогда и только тогда,
когда $\{S_n\}$ ограничена сверху.

\begin{theorem}[признак сравнения]
Пусть $0\le a_k\le b_k$ при $k\ge K_0$. Тогда:
\begin{itemize}[leftmargin=*]
  \item из сходимости $\sum b_k$ следует сходимость $\sum a_k$;
  \item из расходимости $\sum a_k$ следует расходимость $\sum b_k$.
\end{itemize}
\end{theorem}

\begin{proof}
Частичные суммы $A_n\le B_n+C$, где $C$ учитывает первые $K_0$ членов. Если $\sum b_k$ сходится,
$B_n$ ограничена, значит и $A_n$ ограничена — ряд сходится. Второе утверждение — контрапозиция.
\end{proof}

\begin{theorem}[предельный признак сравнения]
Если $a_k,b_k>0$ и $\displaystyle\lim_{k\to\infty}\frac{a_k}{b_k}=K\in(0,+\infty)$, то ряды
$\sum a_k$ и $\sum b_k$ сходятся или расходятся одновременно.
\end{theorem}

\begin{proof}
По $\eps$ найдём $N$: при $k>N$ верно $b_k(K-\eps)<a_k<b_k(K+\eps)$. Применяя признак сравнения к
обоим неравенствам, получаем эквивалентность поведения рядов.
\end{proof}

\begin{theorem}[признак Даламбера]
Пусть $a_k>0$ и $\displaystyle\lim_{k\to\infty}\frac{a_{k+1}}{a_k}=q$. Тогда при $q<1$ ряд
сходится, при $q>1$ расходится; при $q=1$ признак ответа не даёт.
\end{theorem}

\begin{proof}
$q<1$: возьмём $s$ с $q<s<1$; найдётся $N$, что при $k>N$ $\dfrac{a_{k+1}}{a_k}<s$, тогда
$a_{N+m}<s^m a_N$, и $\sum a_k$ мажорируется сходящейся геометрической прогрессией. $q>1$: при
$k>N$ члены возрастают, $a_k\not\to0$, нарушено необходимое условие.
\end{proof}

\begin{example}
$\displaystyle\sum\frac{2^k}{k!}$: $\dfrac{a_{k+1}}{a_k}=\dfrac{2}{k+1}\to0<1$ — сходится.
$\displaystyle\sum\frac{k^k}{k!}$: $\dfrac{a_{k+1}}{a_k}=\bigl(1+\tfrac1k\bigr)^k\to e>1$ —
расходится. Для $\sum\frac1{k^2}$ отношение $\to1$ — признак бессилен.
\end{example}

\begin{theorem}[радикальный признак Коши]
Пусть $a_k\ge0$ и $\displaystyle\lim_{k\to\infty}\sqrt[k]{a_k}=q$. Тогда при $q<1$ ряд сходится,
при $q>1$ расходится; при $q=1$ ответа нет.
\end{theorem}

\begin{proof}
$q<1$: при $q<s<1$ найдётся $N$, что $\sqrt[k]{a_k}<s$, т.е. $a_k<s^k$ — мажорирование
геометрической прогрессией. $q>1$: $a_k>1$ бесконечно часто, $a_k\not\to0$.
\end{proof}

\begin{example}
$\displaystyle\sum\Bigl(\frac{k}{2k+1}\Bigr)^k$: $\sqrt[k]{a_k}=\dfrac{k}{2k+1}\to\dfrac12<1$ —
сходится. (Радикальный признак сильнее Даламбера: работает и для рядов с колеблющимся
отношением, например $\sum\bigl(\tfrac{1+(-1)^k}{3}\bigr)^k$.)
\end{example}

\begin{theorem}[интегральный признак Коши]
Пусть $f(x)$ положительна, непрерывна и монотонно убывает при $x\ge1$, и $a_k=f(k)$. Тогда ряд
$\displaystyle\sum_{k=1}^\infty f(k)$ и интеграл $\displaystyle\int_1^{+\infty}f(x)\dx$ сходятся
или расходятся одновременно.
\end{theorem}

\begin{proof}
В силу монотонности на $[k,k+1]$: $f(k+1)\le\int_k^{k+1}f\le f(k)$. Суммируя,
\[
  S_{n+1}-f(1)\le\int_1^{n+1}f(x)\dx\le S_n ,
\]
откуда $f(n+1)+\int_1^{n+1}f\le S_{n+1}\le f(1)+\int_1^{n+1}f$. Частичные суммы и интеграл
ограничены или неограничены одновременно.
\end{proof}

\subsection{Обобщённый гармонический ряд (ряд Дирихле)}

\begin{theorem}
Ряд $\displaystyle\sum_{k=1}^\infty\frac1{k^p}$ сходится тогда и только тогда, когда $p>1$.
\end{theorem}

\begin{proof}
Применим интегральный признак к $f(x)=x^{-p}$: ряд ведёт себя как
$\int_1^{+\infty}\dfrac{dx}{x^p}$, который сходится при $p>1$ и расходится при $p\le1$ (эталонный
интеграл первого рода).
\end{proof}

\end{document}
